\documentclass[12pt]{article}
\usepackage[a4paper, margin=2.54cm]{geometry}
\usepackage{times}
\usepackage[T1]{fontenc}
\usepackage[utf8]{inputenc}
\usepackage{amsmath}
\usepackage{hyperref}
\usepackage{setspace}
\hypersetup{colorlinks=true,urlcolor=blue}
\setlength{\parskip}{6pt}
\setlength{\parindent}{0pt}
\pagestyle{empty}

\begin{document}

\begin{flushright}
Dvir Ross\\
Department of Software Engineering\\
Shenkar College of Engineering, Design and Art\\
Ramat Gan, Israel\\[2pt]
Department of Computer Science\\
SCE --- Shamoon College of Engineering\\
Beer Sheva, Israel\\[2pt]
\texttt{dvirross@shenkar.ac.il}\\
\texttt{RossDv@sce.ac.il}\\[6pt]
\today
\end{flushright}

\bigskip

The Editors\\
\textit{Statistical Methods in Medical Research}

\bigskip

Dear Editors,

\medskip

I am pleased to submit the manuscript \textbf{``Statistical Evidence for Halachic Vesset
Patterns in Natural Menstrual Cycle Data: A Study of Jewish Religious Law and Menstrual
Cycle Regularity --- A Permutation and Multinomial Sampling Analysis''} for consideration
for publication in \textit{Statistical Methods in Medical Research}.

\medskip
\textbf{What this paper does.}
Jewish religious law (\textit{Halacha}) prescribes that women anticipate menstruation
based on recurring patterns in their cycle history, collectively known as \textit{vesset}.
Despite their centuries-long practical application, no published study has formally tested
whether these patterns occur in real menstrual cycle data at rates exceeding chance.
This paper fills that gap using three complementary null models --- a global permutation
null ($H_G$), a multinomial sampling null ($H_{iid}$), and a within-woman permutation
null ($H_W$) that conditions on each woman's observed cycle-length multiset --- applied
to a publicly available dataset of 1,554 cycles from 118 women
(Fehring et al., \textit{Contraception}, 2013).

\medskip
\textbf{Methodological contributions.}
The paper makes three methodological contributions that we believe are of broad interest
to readers of this journal:

\begin{enumerate}
  \item \textbf{A joint multivariate test for dependent count vectors.}
    We introduce a Mahalanobis-distance-based joint test ($D_M$) that evaluates whether a
    vector of sequentially computed pattern counts departs from the null distribution,
    without assuming independence among the counts. This test is applicable to any setting
    where multiple sequential pattern types are counted in time-series or longitudinal data.

  \item \textbf{A polynomial hierarchy of pattern types.}
    We show that the five vesset pattern types form a hierarchy based on the order of the
    finite-difference operator --- degree-0 (constant), degree-1 (linear), degree-2 (quadratic)
    sequences. We derive three mathematically necessary dependence relationships between
    the patterns (set-theoretic inclusion, limiting convergence, logical incompatibility) and
    confirm each with pairwise chi-square tests. This motivates the joint test and explains
    why Bonferroni correction would be logically inappropriate.

  \item \textbf{A within-woman permutation null as the primary scientific test.}
    We introduce a within-woman null ($H_W$) that permutes each woman's cycle-length
    sequence within her own record, conditioning on her observed cycle-length multiset.
    $H_W$ is the most scientifically conservative of the three nulls: it controls for
    each woman's personal cycle-length distribution and tests only whether patterns arise
    more often than chance within her own sequence. Under $H_W$, no pattern is individually
    significant (all $p \geq .722$) and the joint Mahalanobis test is non-significant
    ($D_M = 2.95$, $p = .121$). A supplementary multinomial null ($H_{iid}$) draws each
    woman's cycle sequence from the fitted empirical distribution; the close agreement of
    $H_G$ and $H_{iid}$ across all five patterns provides a robustness check on the global
    permutation results.
\end{enumerate}

\medskip
\textbf{Main findings.}
The primary result is that no vesset pattern is individually or jointly significant under
the within-woman null ($H_W$): all individual $p \geq .722$ and the joint Mahalanobis
test yields $D_M = 2.95$, $p = .121$. This shows that the observed pattern counts are
consistent with within-woman randomness once each woman's personal cycle-length distribution
is controlled for.

Under the global permutation null ($H_G$), Haflaga departs significantly from chance
($z = 4.51$, $p < .001$; only 2 of 50,000 simulations reaching the observed count).
Dilug is nominally significant ($p = .043$) but does not survive Holm correction
($p_{\text{Holm}} = .173$) and should be treated as hypothesis-generating only.
A focused 3-pattern Mahalanobis test detects significant multivariate departure from
the global null ($D_M = 4.52$, $p < .001$). The contrast between $H_G$ and $H_W$
is itself informative: it shows that Haflaga overrepresentation relative to the population
pool reflects partly the composition of individual women's cycle-length distributions,
rather than purely within-sequence patterning.

\medskip
\textbf{Fit with \textit{Statistical Methods in Medical Research}.}
This paper sits squarely within the scope of the journal: it develops and applies
novel statistical methodology --- randomization testing, multivariate distance tests,
and null model comparison --- to a medical dataset with direct clinical and public-health
relevance (menstrual cycle regularity). The interdisciplinary context (Jewish religious
law) is unusual, but the statistical methodology and its contributions are entirely
general and applicable beyond the halachic setting.

\medskip
\textbf{Originality and prior publication.}
This manuscript has not been submitted elsewhere and is not under review at any other
journal. A preliminary analysis appears in the author's PhD dissertation (Ariel University,
2022, in Hebrew), but the present paper extends that work substantially: it introduces the
joint multivariate test, the multinomial null model, the polynomial hierarchy with
mathematical derivations, the chi-square dependence analysis, and the 5-dimensional
geometric visualization, none of which appeared in the dissertation.

\medskip
\textbf{Data and code availability.}
All analysis code and the filtered dataset are publicly available at
\url{https://github.com/dvirross/VessetVectorTest}.
The original dataset is archived at Marquette University
(\url{https://epublications.marquette.edu/data_nfp/7/}).

\medskip
\textbf{Suggested reviewers.}
Given the interdisciplinary nature of this work, I suggest reviewers with expertise in
both statistical methodology and reproductive medicine or longitudinal biological data.
I am happy to suggest specific names if the editors find this helpful.

\medskip

I believe this manuscript will be of genuine interest to the readership of
\textit{Statistical Methods in Medical Research}, and I look forward to the
editorial team's consideration.

\bigskip

Yours sincerely,

\bigskip\bigskip

Dvir Ross, PhD\\
Department of Software Engineering, Shenkar College of Engineering, Design and Art\\
Department of Computer Science, SCE --- Shamoon College of Engineering

\newpage

% ── DECLARATIONS PAGE ─────────────────────────────────────────────────────────
\begin{center}
  {\large\bfseries Declarations}
\end{center}

\bigskip

\textbf{Funding}

This research received no specific grant from any funding agency in the public,
commercial, or not-for-profit sectors.

\bigskip

\textbf{Conflict of interest}

The author declares no conflicts of interest.

\bigskip

\textbf{Ethics approval}

This study is a secondary analysis of a publicly available, fully anonymised dataset
(Fehring et al., 2013, archived at Marquette University). No new data collection
was conducted. No ethical approval was required.

\bigskip

\textbf{Patient consent}

Not applicable. No new data were collected. The original dataset was collected under
the ethical oversight of the original investigators (Fehring et al., 2013).

\bigskip

\textbf{Data availability}

The filtered dataset and all analysis code used in this study are publicly available
at \url{https://github.com/dvirross/VessetVectorTest}. The original NFP dataset is
archived at Marquette University:
\url{https://epublications.marquette.edu/data_nfp/7/}.

\bigskip

\textbf{Author contributions}

Dvir Ross: Conceptualization, Methodology, Software, Formal analysis, Investigation,
Writing --- original draft, Writing --- review \& editing, Visualization.

\bigskip

\textbf{Artificial intelligence use disclosure}

Portions of the written text in this article were drafted and edited with the assistance
of Claude (Anthropic, 2024), a large language model AI assistant. All scientific
content, mathematical derivations, statistical analyses, and intellectual contributions
are solely the work of the author. The AI tool was used exclusively for language
editing, structural organisation of the manuscript, and generation of \LaTeX{} and
Python code under the author's direct supervision and verification. The author takes
full responsibility for the accuracy and integrity of all content presented in this
article.

\bigskip

\textbf{ORCID}

Dvir Ross: \href{https://orcid.org/0000-0001-8217-3359}{0000-0001-8217-3359}

\end{document}
